% Part: computability
% Chapter: recursive-functions
% Section: pr-functions

\documentclass[../../../include/open-logic-section]{subfiles}

\begin{document}

\olfileid[fr]{cmp}{rec}{prf}
\olsection{Fonctions récursives primitives}

Récapitulons les deux procédés qui permettent de définir de nouvelles
fonctions à partir de fonctions connues : la récurrence primitive
et la composition.

\begin{defn}
  \ollabel{defn:primitive-recursion} Soit $f$ une fonction à $k$
  arguments ($k\ge 1$) et $g$ une fonction à $k+2$ arguments.
  La fonction \emph{définie par récurrence primitive à partir de
  $f$ et de $g$} est la fonction~$h$ à $k+1$ arguments donnée par
  les équations
  \begin{align*}
    h(x_0,\dots,x_{k-1},0) & =  f(x_0,\dots,x_{k-1}) \\
    h(x_0,\dots,x_{k-1},y+1) & = g(x_0,\dots,x_{k-1}, y, h(x_0,\dots,x_{k-1}, y))
  \end{align*}
\end{defn}

\begin{defn}
  \ollabel{defn:composition} Soit $f$ une fonction à $k$ arguments
  et soient $g_0,\dots,g_{k-1}$ des fonctions à $n$ arguments.
  La fonction \emph{définie par composition de $f$ avec
    $g_0,\dots,g_{k-1}$} est la fonction~$h$ à $n$ arguments
  donnée par
  \[
  h(x_0, \dots, x_{n-1}) =
  f(g_0(x_0, \dots, x_{n-1}), \dots, g_{k-1}(x_0, \dots, x_{n-1})).
  \]
\end{defn}

Outre la fonction successeur~$\Succ$ et les fonctions de projection
\[
\Proj{n}{i}(x_0,\dots,x_{n-1}) = x_i,
\]
pour chaque entier naturel~$n$ et chaque $i<n$, nous compterons
parmi les fonctions récursives primitives la fonction nulle
$\Zero(x)=0$.

\begin{defn}
  L'ensemble des fonctions récursives primitives est l'ensemble
  des fonctions de $\Nat^n$ dans $\Nat$, pour $n\geq1$, défini inductivement
  par les clauses suivantes :
  \begin{enumerate}
  \item $\Zero$ est récursive primitive.
  \item $\Succ$ est récursive primitive.
  \item Chaque projection $\Proj{n}{i}$ est récursive primitive.
  \item Si $f$ est une fonction récursive primitive à $k$ arguments
    et si $g_0,\dots,g_{k-1}$ sont des fonctions récursives
    primitives à $n$ arguments, la composition de $f$ avec
    $g_0,\dots,g_{k-1}$ est récursive primitive.
  \item Si $f$ est une fonction récursive primitive à $k$ arguments
    et $g$ une fonction récursive primitive à $k+2$ arguments,
    la fonction définie par récurrence primitive à partir de
    $f$ et de $g$ est récursive primitive.
  \end{enumerate}
\end{defn}

\begin{explain}
En bref, l'ensemble des fonctions récursives primitives est le
plus petit ensemble qui contient $\Zero$, $\Succ$ et les
projections~$\Proj{n}{j}$, et qui est clos par composition et
par récurrence primitive.

On peut aussi décrire cet ensemble par « étapes ». Soit $S_0$
l'ensemble des fonctions de base : $\Zero$, $\Succ$ et les
projections. Ce sont les fonctions récursives primitives de
l'étape~$0$. Une fois $S_i$ défini, on obtient $S_{i+1}$ en
conservant les fonctions de~$S_i$ et en y ajoutant toutes celles
que donne une seule application de la composition ou de la
récurrence primitive à des fonctions de~$S_i$. Alors
\[
S = \bigcup_{i \in \Nat} S_i
\]
est l'ensemble de toutes les fonctions récursives primitives.\footnote{Explicitation éditoriale : les étapes décrites dans la source sont déjà croissantes, car la composition de la fonction identité avec une fonction la reproduit, et l'identité reste présente à chaque étape. Nous rendons cette conservation explicite ; nous ne changeons pas la classe obtenue.}
\end{explain}

Vérifions que $\Add$ est récursive primitive.

\begin{prop}
  La fonction d'addition $\Add(x,y)=x+y$ est récursive primitive.
\end{prop}

\begin{proof}
Nous disposons déjà d'une définition par récurrence primitive
de $\Add$ au moyen de deux fonctions $f$ et~$g$; elle a la
forme exigée par la \olref{defn:primitive-recursion} :
\begin{align*}
  \Add(x_0, 0) & = f(x_0) = x_0\\
  \Add(x_0, y+1) & = g(x_0, y, \Add(x_0, y)) =
  \Succ(\Add(x_0, y))
\end{align*}
Il suffit donc de montrer que $f$ et $g$ sont elles-mêmes
récursives primitives. Or $f(x_0)=x_0=\Proj{1}{0}(x_0)$ :
les projections étant récursives primitives, $f$ l'est.
La fonction~$g$ à trois arguments est définie par
\[
g(x_0, y, z) = \Succ(z).
\]
Cela ne suffit pas encore à conclure : $\Succ$ n'a qu'un
argument, tandis que $g$ en a trois. Mais on peut définir
$g$ par composition, de façon explicite :
\[
g(x_0, y, z) = \Succ(\Proj{3}{2}(x_0, y, z))
\]
Puisque $\Succ$ et $\Proj{3}{2}$ sont récursives primitives,
leur composition~$g$ l'est aussi.
\end{proof}

\begin{prop}
  \ollabel{prop:mult-pr}
  La fonction de multiplication $\Mult(x,y)=x\cdot y$ est
  récursive primitive.
\end{prop}

\begin{proof}
  Exercice.
\end{proof}

\begin{prob}
  Démontrer la \olref[cmp][rec][prf]{prop:mult-pr} en mettant la
  définition par récurrence primitive de $\Mult$ sous la forme
  exigée par la \olref[cmp][rec][prf]{defn:primitive-recursion},
  puis en montrant que les fonctions $f$ et~$g$ correspondantes
  sont récursives primitives.
\end{prob}

\begin{ex}
Reprenons notre tout premier exemple de définition récursive
primitive :
\begin{align*}
h(0) & =  1 \\
h(y+1) & =  2 \cdot h(y).
\end{align*}
Cette fonction n'entre pas dans le schéma de la
\olref{defn:primitive-recursion}, car $k=0$. Sa définition
fait aussi intervenir les constantes $1$ et~$2$. Pour surmonter
la première difficulté, ajoutons un argument fictif et
définissons~$h'$ :
\begin{align*}
h'(x_0, 0) & =  f(x_0) = 1 \\
h'(x_0, y+1) & =  g(x_0, y, h'(x_0, y)) = 2 \cdot h'(x_0, y).
\end{align*}
La fonction $f(x_0)=1$ s'obtient par composition de $\Succ$
et de $\Zero$ : $f(x_0)=\Succ(\Zero(x_0))$. La fonction~$g$
s'obtient par composition à partir de $g'(z)=2\cdot z$ et
des projections :
\begin{align*}
g(x_0, y, z) & = g'(\Proj{3}{2}(x_0, y, z))
\intertext{et $g'$ s'obtient à son tour par composition :}
g'(z) & = \Mult(g''(z), \Proj{1}{0}(z))
\intertext{où}
g''(z) & = \Succ(f(z)),
\end{align*}
$f$ étant la fonction définie ci-dessus :
$f(z)=\Succ(\Zero(z))$. Une fois $h'$ obtenue, la composition
permet encore de poser
$h(y)=h'(\Proj{1}{0}(y),\Proj{1}{0}(y))$.
Ainsi, $h$ se définit à partir des fonctions de base par une
suite de compositions et de récurrences primitives : elle est
donc récursive primitive.
\end{ex}

\end{document}
